\documentclass[11pt,a4paper]{article}
\usepackage[utf8]{inputenc}
\usepackage{amsmath, amssymb, amsthm}
\usepackage{graphicx}
\usepackage{subcaption}
\usepackage{enumitem}
\usepackage{xcolor}
\usepackage{colortbl}
\usepackage{hyperref}

\newtheorem{proposition}{Proposition}[section]

\title{THE SOLAR SYSTEM}
\author{Johannes Kepler\footnote{Imperial Mathematician, Prague}}
\date{}

\begin{document}

\maketitle
\tableofcontents
\newpage

\section{Our Planetary Neighborhood}
\subsection{The Sun and Its Planets}
The Sun and its gravity-bound bodies form the Solar System \cite{kepler}. See \href{https://science.nasa.gov/solar-system/}{NASA Solar System Exploration}. In increasing distance, the planets are
Mercury, Venus, Earth, Mars, Jupiter, Saturn, Uranus, Neptune.

The planets form two groups:
\begin{enumerate}[label=(\Roman*)]
    \item Inner planets
    \begin{itemize}[label=$\dagger$]
        \item Mercury, Venus, Earth, and Mars belong to this group.
        \item They are comparatively small and have solid surfaces.
        \item They are also called the terrestrial planets.
    \end{itemize}
    \item Outer planets
    \begin{itemize}[label=$\dagger$]
        \item Jupiter and Saturn are gas giants.
        \item Uranus and Neptune are commonly called ice giants.
        \item All four outer planets possess ring systems.
    \end{itemize}
\end{enumerate}
The inner planets are small, dense, and relatively close to the Sun, whereas the outer planets are much larger and travel along considerably wider orbits.

\subsection{Portraits of the Solar System}
Figure \ref{fig:solarsystem} shows the Sun and eight planets.\footnote{Sizes and distances are schematic.} Compare Mars in Figure \ref{fig:mars} with Jupiter in Figure \ref{fig:jupiter}.

\section{Kepler's Laws of Planetary Motion}
Kepler's three laws describe planetary motion \cite{nasa_orbits}.

\subsection{First Law: Elliptical Orbits}
A planetary orbit is elliptical, with the Sun at one focus:
\begin{equation}
\frac{x^{2}}{a^{2}}+\frac{y^{2}}{b^{2}}=1, \quad a\ge b>0
\end{equation}
Here $a$ and $b$ are the semi-axes. 

\begin{figure}[htbp]
    \centering
    \textbf{\Large Star}\par\medskip
    \begin{subfigure}{0.3\textwidth}
        \centering
        \includegraphics[width=\linewidth]{sun.png}
        \caption{The Sun}
    \end{subfigure}
    
    \vspace{1cm}
    \textbf{\Large Inner Planets}\par\medskip
    \begin{subfigure}{0.2\textwidth}
        \centering
        \includegraphics[width=\linewidth]{mercury.png}
        \caption{Mercury}
    \end{subfigure}\hfill
    \begin{subfigure}{0.2\textwidth}
        \centering
        \includegraphics[width=\linewidth]{venus.png}
        \caption{Venus}
    \end{subfigure}\hfill
    \begin{subfigure}{0.2\textwidth}
        \centering
        \includegraphics[width=\linewidth]{earth.png}
        \caption{Earth}
    \end{subfigure}\hfill
    \begin{subfigure}{0.2\textwidth}
        \centering
        \includegraphics[width=\linewidth]{mars.png}
        \caption{Mars}
        \label{fig:mars}
    \end{subfigure}
    
    \vspace{1cm}
    \textbf{\Large Outer Planets}\par\medskip
    \begin{subfigure}{0.2\textwidth}
        \centering
        \includegraphics[width=\linewidth]{jupiter.png}
        \caption{Jupiter}
        \label{fig:jupiter}
    \end{subfigure}\hfill
    \begin{subfigure}{0.2\textwidth}
        \centering
        \includegraphics[width=\linewidth]{saturn.png}
        \caption{Saturn}
    \end{subfigure}\hfill
    \begin{subfigure}{0.2\textwidth}
        \centering
        \includegraphics[width=\linewidth]{uranus.png}
        \caption{Uranus}
    \end{subfigure}\hfill
    \begin{subfigure}{0.2\textwidth}
        \centering
        \includegraphics[width=\linewidth]{neptune.png}
        \caption{Neptune}
    \end{subfigure}
    
    \caption{The Sun and the eight planets of the Solar System, arranged into the inner and outer planetary groups. The displayed sizes and distances are not to scale.}
    \label{fig:solarsystem}
\end{figure}

For focal distance $c$,
\begin{equation}
c^{2}=a^{2}-b^{2}, \quad e=\frac{c}{a}
\end{equation}
Here $e$ is the eccentricity. In polar form,
\begin{equation}
r(\theta)=\frac{a(1-e^{2})}{1+e\cos\theta}.
\end{equation}
Equation (3) gives the orbital distance.

\subsection{Second Law: Equal Areas in Equal Times}
Equal times sweep equal areas:
\begin{equation}
\Delta t_1 = \Delta t_2 \implies \Delta A_1 = \Delta A_2.
\end{equation}
Consequently, a planet moves
\begin{itemize}
    \item faster when it is closer to the Sun, and
    \item slower when it is farther from the Sun.
\end{itemize}

\subsection{Third Law: Orbital Period and Distance}
For planets orbiting one star, $T^{2}\propto a^{3}$. Newtonian form:
\begin{equation}
T^{2}=\frac{4\pi^{2}}{GM_{\odot}}a^{3}
\end{equation}
Here $M_{\odot}$ is the solar mass.

\subsection{An Earth-Mars Comparison}
Using Earth as reference, $a_{E}=1\text{ AU}$, $T_{E}=1\text{ year}$, where $1\text{ AU}\approx1.496\times10^{11}\text{ m}$.\footnote{AU denotes the mean Earth-Sun distance.} For Mars, $a_{M}\approx1.524\text{ AU}$, so
\begin{equation}
\frac{T_{M}^{2}}{T_{E}^{2}}=\frac{a_{M}^{3}}{a_{E}^{3}}.
\end{equation}
Thus $T_{M}=T_{E}\left(\frac{a_{M}}{a_{E}}\right)^{3/2} \approx (1.524)^{3/2} \approx 1.88\text{ years}$.

\section{A Relativistic Extension}
General relativity extends Kepler's model.\footnote{Planetary corrections are small but measurable.} Einstein's equations are
\begin{equation}
G_{\mu\nu}+\Lambda g_{\mu\nu}=\frac{8\pi G}{c^{4}}T_{\mu\nu}
\end{equation}
For spherical mass $M$,
\begin{equation}
ds^{2}=-\left(1-\frac{r_{s}}{r}\right)c^{2}dt^{2}+\left(1-\frac{r_{s}}{r}\right)^{-1}dr^{2}+r^{2}(d\theta^{2}+\sin^{2}\theta d\phi^{2}), \quad r_{s}=\frac{2GM}{c^{2}}.
\end{equation}
Free fall satisfies
\begin{equation}
\frac{d^{2}x^{\mu}}{d\tau^{2}}+\Gamma_{\alpha\beta}^{\mu}\frac{dx^{\alpha}}{d\tau}\frac{dx^{\beta}}{d\tau}=0.
\end{equation}
Per orbit,
\begin{equation}
\Delta\varpi=\frac{6\pi GM}{a(1-e^{2})c^{2}}
\end{equation}

\begin{proposition}[Weak-field circular orbit]
For a circular orbit of radius $r$,
\begin{equation}
\left(\frac{v}{c}\right)^{2}=\frac{GM}{rc^{2}}=\frac{r_{s}}{2r}.
\end{equation}
\end{proposition}
\begin{proof}
The circular-orbit balance $v^{2}/r=GM/r^{2}$ gives $v^{2}=GM/r$. Divide by $c^{2}$ and use $r_{s}=2GM/c^{2}$.
\end{proof}
See Proposition 3.1.

\section{The Outer Planets}
The outer planets are Jupiter, Saturn, Uranus, and Neptune \cite{nasa_planets}. The first two are gas giants; the last two are ice giants.

Table \ref{tab:outer_planets} uses scientific and relativistic notation.\footnote{Values are rounded for typesetting practice.}

\begin{table}[h!]
\centering
\caption{Orbital and relativistic parameters of the outer planets}
\label{tab:outer_planets}
\begin{tabular}{|l|c|c|c|c|}
\hline
\rowcolor{blue!15}
\textbf{Planet} & \multicolumn{2}{c|}{\textbf{Orbital parameters}} & \multicolumn{2}{c|}{\textbf{Relativistic notation}} \\ \cline{2-5}
\rowcolor{blue!5}
 & $a$ (AU) & $e$ & $\epsilon=\frac{1}{ac^{2}}GM_{\odot}$ & $\beta=\sqrt{\epsilon}$ \\ \hline
Jupiter & $5.204$ & $4.89\times10^{-2}$ & $1.90\times10^{-9}$ & $4.36\times10^{-5}$ \\ \hline
Saturn & $9.583$ & $5.65\times10^{-2}$ & $1.03\times10^{-9}$ & $3.21\times10^{-5}$ \\ \hline
Uranus & $19.191$ & $4.72\times10^{-2}$ & $5.14\times10^{-10}$ & $2.27\times10^{-5}$ \\ \hline
Neptune & $30.070$ & $8.6\times10^{-3}$ & $3.28\times10^{-10}$ & $1.81\times10^{-5}$ \\ \hline
\end{tabular}
\end{table}

Table \ref{tab:outer_planets} compares the four giant planets.

\begin{thebibliography}{9}
\bibitem{kepler}
J. Kepler, \textit{Astronomia Nova}. Heidelberg, 1609.
\bibitem{nasa_orbits}
National Aeronautics and Space Administration, "Orbits and kepler's laws," 2024.
\bibitem{nasa_planets}
National Aeronautics and Space Administration, "About the planets," 2026. Solar System Exploration.
\end{thebibliography}

\end{document}